\chapter{实现要点}
\label{ch:impl}

\section{技术栈与运行形态}
后端采用 C++20，基于 HTTP 暴露初始化、插入、查询、删除与统计接口（依赖如 \texttt{cpp-httplib}\cite{cpphttplib} 与 \texttt{nlohmann/json}\cite{nlohmannjson}）；持久化与实验数据使用 MySQL\cite{mysqlmanual}；前端采用 Vue 3\cite{vue3} 组织实验面板与图表展示。该划分便于将“热路径内存结构”与“冷路径审计落库”解耦。

\section{核心数据结构映射}
第 \ref{ch:design} 章所述类型与实现一一对应：全局 \texttt{HashTable} 持有 \(n,N,K\)、\texttt{Facility*} 数组、\texttt{ResizeManager}、\texttt{RebuildScheduler} 及可逆置换 \texttt{PermutationHash}（32 位分半 Feistel 三轮，密钥字段 \texttt{k1..k4}）。每个 \texttt{Facility} 含 \texttt{MiniArray D}、\texttt{cubbies} 列表、\texttt{tail} 指针及 \texttt{tier\_counts}。\texttt{Cubby} 含 \texttt{tier}、\texttt{capacity}、\texttt{size}、载荷向量 \texttt{S}、\texttt{MiniArray Meta} 与 \texttt{FreeSlotTree}（线性化位图，支持 \texttt{find\_free}/\texttt{mark\_used}/\texttt{mark\_free}）。

\texttt{MetaEntry} 逻辑上划分为商压缩字段（\texttt{delta}、\texttt{mid\_bits}）、路由辅助（\texttt{router\_bits}）与插入元数据（\texttt{insert\_bits}）；编码为变长比特串后写入 MiniArray。MiniArray 内部为 packed B-树加查表，满足“总空间与总比特数线性、单点访问均摊常数”的工程目标。

\section{算法伪代码与代码对应}
插入路径与第 \ref{ch:design} 章一致：计算 \(g_x=\pi(x)\)，拆分设施索引 \(r\) 与路由索引 \(b\)，取 \texttt{F.D[b]}，在 tail cubby 内完成 k-kick 探测循环，写 \texttt{S[slot]} 与 \texttt{Meta}，更新路由器，调用不变量维护。查询路径仅经路由器定位 \((I,i)\)，再 \texttt{Meta.access(i)} 与 \texttt{reconstruct}。删除路径在清槽后按是否 tail 决定是否从 tail 搬迁，并在 tail 空时回收与晋升新 tail。

\section{不变量维护与重建}
\texttt{maintainInvariants(Facility*)} 遍历 tier，若计数低于 \(t_j/2\) 调用 \texttt{rebuildUp(j)}（拆分高 tier），高于 \(3t_j\) 调用 \texttt{rebuildDown(j)}（合并一组 \(j\) 级 cubby）。实现细节包括：合并时创建 \(j{+}1\) 级新 cubby 并逐元素插入；拆分时从高 tier 取一枚大 cubby，随机分配元素到 \(t_j\) 个小 cubby 并分别重建。

\section{扩容与缩容}
当 \(n\) 超过当前容量上界触发向新表迁移（\texttt{ResizeManager} 持有 \texttt{newh} 与进度）；当 \(n<\) 容量 \(/4\) 触发收缩。双表迁移将重建与前台操作交错，避免长时间停顿（具体调度以实现为准）。

\section{数据库与观测}
见第 \ref{ch:exp} 章与附录：指标表 \texttt{metrics} 记录每次操作的 \texttt{probe\_count}、\texttt{kick\_count}、延迟与 \texttt{cubby\_tier}；快照表记录 cubby、tier 与槽位占用，用于验证分布与负载。写入采用线程本地缓冲批量 flush，避免每次操作直连数据库。
